\documentclass[11pt,a4paper]{article}
\usepackage[utf8]{inputenc}
\usepackage[vietnamese]{babel}
\usepackage{amsmath,amssymb}
\usepackage{geometry}
\geometry{a4paper, total={180mm,260mm}, left=15mm, right=15mm, top=18mm, bottom=18mm}

\usepackage{tikz}
\usetikzlibrary{shapes.geometric, arrows.meta, positioning, calc}
\tikzset{arr/.style={-{Stealth[scale=0.8]}}}
\usepackage{tcolorbox}
\tcbuselibrary{skins,breakable}
\usepackage{fontawesome}
\usepackage{fancyhdr}
\usepackage{enumitem}

\pagestyle{fancy}
\fancyhf{}
\renewcommand{\headrulewidth}{0.5pt}
\renewcommand{\footrulewidth}{0.5pt}
\fancyhead[L]{\small\bfseries\color{violet!85!black}\faBook\ SỔ TAY KIẾN THỨC TOÁN 2}
\fancyhead[R]{\small\bfseries\color{violet!85!black}Lớp Toán Thầy Lĩnh -- loptoanthaylinh.vn}
\fancyfoot[C]{\small\color{gray}-- Sổ tay kiến thức: Bài 10. Bảng cộng có nhớ trong phạm vi 20 --}

\setlength{\parindent}{0pt}
\setlength{\parskip}{0pt}

\begin{document}

% --- BANNER BÀI HỌC ---
\begin{tcolorbox}[
  enhanced,
  colback=violet!8!white,
  colframe=violet!70!black,
  arc=3mm, boxrule=1.2pt,
  left=3mm, right=3mm, top=2mm, bottom=2mm,
  halign=center
]
  {\LARGE\bfseries\color{violet!90!black}\faBookmark\ BÀI 10: BẢNG CỘNG (CÓ NHỚ) TRONG PHẠM VI 20}\\[1mm]
  {\itshape\footnotesize\color{black!75} Hệ thống 4 bảng cộng cửu chương qua 10 và kĩ năng tính nhẩm nhanh trên 1 trang A4}
\end{tcolorbox}

\vspace{-2mm}

\noindent
\begin{minipage}[t]{0.49\textwidth}\vspace{0pt}
\begin{tcolorbox}[
  enhanced,
  colback=violet!4!white,
  colframe=violet!75!black,
  arc=2mm, boxrule=0.8pt,
  left=2mm, right=2mm, top=1.8mm, bottom=1.8mm,
  title={\footnotesize\bfseries\color{white}\faKey\ 1. Bảng 9 cộng và Bảng 8 cộng}
]
{\footnotesize
\begin{tabular}{l|l}
  \textbf{Bảng 9 cộng:} & \textbf{Bảng 8 cộng:} \\ \hline
  $9 + 2 = 11$ & $8 + 3 = 11$ \\
  $9 + 3 = 12$ & $8 + 4 = 12$ \\
  $9 + 4 = 13$ & $8 + 5 = 13$ \\
  $9 + 5 = 14$ & $8 + 6 = 14$ \\
  $9 + 6 = 15$ & $8 + 7 = 15$ \\
  $9 + 7 = 16$ & $8 + 8 = 16$ \\
  $9 + 8 = 17$ & $8 + 9 = 17$ \\
  $9 + 9 = 18$ &
\end{tabular}
}
\end{tcolorbox}
\end{minipage}
\hfill
\begin{minipage}[t]{0.49\textwidth}\vspace{0pt}
\begin{tcolorbox}[
  enhanced,
  colback=teal!4!white,
  colframe=teal!75!black,
  arc=2mm, boxrule=0.8pt,
  left=2mm, right=2mm, top=1.8mm, bottom=1.8mm,
  title={\footnotesize\bfseries\color{white}\faSliders\ 2. Bảng 7 cộng và Bảng 6 cộng}
]
{\footnotesize
\begin{tabular}{l|l}
  \textbf{Bảng 7 cộng:} & \textbf{Bảng 6 cộng:} \\ \hline
  $7 + 4 = 11$ & $6 + 5 = 11$ \\
  $7 + 5 = 12$ & $6 + 6 = 12$ \\
  $7 + 6 = 13$ & $6 + 7 = 13$ \\
  $7 + 7 = 14$ & $6 + 8 = 14$ \\
  $7 + 8 = 15$ & $6 + 9 = 15$ \\
  $7 + 9 = 16$ & \\
  \multicolumn{2}{l}{\textit{*Các bảng khác:}} \\
  $5 + 6 = 11$; $5 + 7 = 12$; $5 + 8 = 13$; $5 + 9 = 14$
\end{tabular}
}
\end{tcolorbox}
\end{minipage}

\vspace{-1.5mm}
\begin{tcolorbox}[
  enhanced,
  colback=orange!4!white,
  colframe=orange!80!black,
  arc=2mm, boxrule=0.8pt,
  left=2mm, right=2mm, top=1.5mm, bottom=1.5mm,
  title={\footnotesize\bfseries\color{white}\faSitemap\ 3. Sơ đồ tư duy: Tính chất giao hoán trong bảng cộng}
]
\centering
\begin{tikzpicture}[x=1.5cm, y=0.8cm]
  \node[rectangle, rounded corners=2pt, draw=violet!80!black, fill=violet!15, font=\scriptsize\bfseries] (a) at (0,0) {$9 + 4 = 13$};
  \node[rectangle, rounded corners=2pt, draw=teal!80!black, fill=teal!15, font=\scriptsize\bfseries] (b) at (3,0) {$4 + 9 = 13$};
  \draw[<->, very thick, orange!80!black] (a) -- node[above, font=\tiny\bfseries] {Đổi chỗ số hạng} node[below, font=\tiny\bfseries] {Tổng không đổi} (b);

  \node[rectangle, rounded corners=2pt, draw=violet!80!black, fill=violet!15, font=\scriptsize\bfseries] (c) at (0,-1) {$8 + 5 = 13$};
  \node[rectangle, rounded corners=2pt, draw=teal!80!black, fill=teal!15, font=\scriptsize\bfseries] (d) at (3,-1) {$5 + 8 = 13$};
  \draw[<->, very thick, orange!80!black] (c) -- node[above, font=\tiny\bfseries] {Đổi chỗ số hạng} node[below, font=\tiny\bfseries] {Tổng không đổi} (d);
\end{tikzpicture}
\end{tcolorbox}

\vspace{-1.5mm}
\noindent
\begin{minipage}[t]{0.49\textwidth}\vspace{0pt}
\begin{tcolorbox}[
  enhanced,
  colback=green!4!white,
  colframe=green!60!black,
  arc=2mm, boxrule=0.8pt,
  left=2mm, right=2mm, top=1.8mm, bottom=1.8mm,
  title={\footnotesize\bfseries\color{white}\faCheckCircle\ 4. Ví dụ thực hành chuẩn mực}
]
{\footnotesize
\textbf{1. Tính nhẩm:}\\
$9 + 7 = \mathbf{16}$; $8 + 6 = \mathbf{14}$; $7 + 8 = \mathbf{15}$.\\
$6 + 8 = \mathbf{14}$; $5 + 9 = \mathbf{14}$.
\textbf{2. Điền số thích hợp vào ô trống:}\\
\begin{itemize}\setlength{\itemsep}{0.5pt}\setlength{\leftskip}{-2.5mm}
  \item $8 + \mathbf{5} = 13$ (vì $13 - 8 = 5$).
  \item $7 + \mathbf{7} = 14$.
  \item $\mathbf{9} + 6 = 15$.
\end{itemize}
}
\end{tcolorbox}
\end{minipage}
\hfill
\begin{minipage}[t]{0.49\textwidth}\vspace{0pt}
\begin{tcolorbox}[
  enhanced,
  colback=red!4!white,
  colframe=red!75!black,
  arc=2mm, boxrule=0.8pt,
  left=2mm, right=2mm, top=1.8mm, bottom=1.8mm,
  title={\footnotesize\bfseries\color{white}\faExclamationTriangle\ 5. Mẹo tránh bẫy trắc nghiệm}
]
{\footnotesize
\begin{itemize}\setlength{\itemsep}{1pt}\setlength{\leftskip}{-2.5mm}
  \item \textbf{Mẹo cộng 3 số: Nhóm số có tổng bằng 10:}\\
  Tính: $8 + 7 + 2 = (8 + 2) + 7 = 10 + 7 = 17$.\\
  Tính: $9 + 5 + 1 = (9 + 1) + 5 = 10 + 5 = 15$.
  \item Khi gặp số bé đứng trước (như $3 + 8$), hãy nhẩm đổi lại thành $8 + 3 = 11$ để tính nhanh hơn.
  \item Học thuộc lòng các bảng cộng để làm nền tảng cho phép cộng có nhớ lớp lớn.
\end{itemize}
}
\end{tcolorbox}
\end{minipage}


\end{document}
